\section{Conclusion}

Embodied intelligence reframes medical AI from a system that interprets observations into one that participates in the physical production of evidence and care. Across monitoring, diagnosis, intervention and rehabilitation, the central opportunity is a closed loop in which sensing, reasoning and action are jointly designed around a clinical need. The same loop also concentrates risk: errors can propagate into the body and environment, continuous observation can erode privacy, and adaptation can outpace conventional validation. Progress will therefore depend less on maximal autonomy than on calibrated agency---systems that know what they can sense, when they should abstain, which actions are permitted and how responsibility returns to people. If grounded in prospective evidence, equitable access, cybersecurity and meaningful oversight, embodied systems could move medicine from episodic response toward continuous, personalized partnership. Without those foundations, physical capability will merely amplify weaknesses already present in data and institutions. The field's task is to make embodiment not only technically competent, but clinically useful, governable and worthy of trust.

